\chapter{Application Pratique~: Prédiction du Rendement Agricole --- Plateforme AgriYield}

\encadreobjectif{Présenter une implémentation complète, cohérente et traçable du thème \textbf{``Exploitation des données à l'aide des modèles d'IA''}, en expliquant clairement l'architecture, les pages de la plateforme, la gouvernance \textit{baseline/fine-tuning} et l'intérêt concret de comparer plusieurs modèles.}

%--------------------------------------------------
\section{Justification du thème et du cas d'usage}

Le thème de base du projet consiste à transformer des données hétérogènes en décisions exploitables. Dans AgriYield, cette logique est matérialisée par trois niveaux complémentaires~:
\begin{itemize}[leftmargin=1.5cm]
  \item \textbf{Niveau données}~: ingestion d'historiques agricoles et consolidation de variables explicatives (pluie, température, pesticides).
  \item \textbf{Niveau IA}~: entraînement multi-modèles, sélection, simulation de scénarios et génération d'alertes.
  \item \textbf{Niveau décision}~: recommandation métier justifiée pour l'action publique/agronomique.
\end{itemize}

Le projet est donc une chaîne décisionnelle complète, et non un simple exercice d'entraînement de modèle.

%--------------------------------------------------
\section{Architecture technique implémentée}

\subsection{Vue microservices}

\begin{figure}[H]
  \centering
  \includegraphics[width=0.95\textwidth]{arcitecture.png}
  \caption{Architecture microservices de la plateforme AgriYield}
  \label{fig:agriyield-architecture-finale}
\end{figure}

\textbf{Lecture de la figure~:} le frontend React passe par un gateway unique qui distribue les requêtes vers les services métier (AgriData, Integration, Decision, ML), avec Discovery/Config pour la gouvernance Spring.

\begin{table}[H]
\centering
\caption{Services, ports et responsabilités}
\label{tab:services-architecture}
\renewcommand{\arraystretch}{1.25}
\begin{tabularx}{\textwidth}{|>{\bfseries}p{2.4cm}|c|>{\RaggedRight\arraybackslash}X|>{\RaggedRight\arraybackslash}X|}
\hline
\rowcolor{TableHeader}
Service & Port & Responsabilité & Endpoints/usage \\
\hline
Frontend React & 5173 & Interface utilisateur (tableaux de bord, simulation, alertes, gouvernance) & --- \\
\hline
Gateway & 8081 & Entrée API unique, routage, CORS, contexte de requête & /api/agri-data/*, /api/integration/*, /api/decision/*, /api/ml/* \\
\hline
ML Python (FastAPI) & 8000 & Entraînement, prédiction, simulateur, alertes, performances, registre & /train/upload, /overview, /scenario/simulate, /alerts, /performance \\
\hline
AgriData & 8093 & Données agricoles structurées et couverture temporelle & /records, /records/coverage, /records/datasets \\
\hline
Integration & 8095 & Signaux externes (météo, marché, NDVI) avec fallback & /signals/realtime \\
\hline
Decision & 8094 & Fusion ML + signaux externes + score de risque + recommandations & /analysis \\
\hline
Discovery Server & 8090 & Registre de services (Eureka) & /eureka \\
\hline
Config Server & 8888 & Configuration centralisée (config-repo) & /\{service\}/\{profile\} \\
\hline
\end{tabularx}
\end{table}

\subsection{Flux décisionnel}

\begin{enumerate}[leftmargin=1.5cm]
  \item L'utilisateur agit sur l'interface React.
  \item Le Gateway route vers le bon service.
  \item Le service ML produit les prédictions et métriques du modèle actif.
  \item Le service Integration fournit les signaux externes avec fraîcheur et confiance.
  \item Le service Decision consolide les contributions et calcule le niveau de risque.
  \item Le frontend restitue les résultats avec contexte de confiance (version, source, couverture temporelle).
\end{enumerate}

%--------------------------------------------------
\section{Description claire de chaque page de la plateforme}

\subsection{Accueil + gouvernance latérale}

\begin{figure}[H]
  \centering
  \includegraphics[width=0.95\textwidth]{accueil.png}
  \caption{Page d'accueil, modèle actif et panneau de confiance}
  \label{fig:agriyield-vue-accueil}
\end{figure}

\textbf{Rôle de cette page~:}
\begin{itemize}[leftmargin=1.5cm]
  \item charger un dataset CSV,
  \item afficher le modèle actif (algorithme, version, précision),
  \item exposer les éléments de traçabilité (source dataset, hash, horizon temporel),
  \item fournir des actions de gouvernance modèle (activation, rollback baseline).
\end{itemize}

\subsection{Page ``Simulateur''}

\begin{figure}[H]
  \centering
  \includegraphics[width=0.95\textwidth]{page_simulation.png}
  \caption{Simulation de scénarios climatiques et analyse de sensibilité}
  \label{fig:agriyield-vue-simulation}
\end{figure}

\textbf{Entrées~:} pays, culture, année cible, variations pluie/température/pesticides.\\
\textbf{Sorties~:} comparaison historique vs base vs scénario, différentiel net, courbes de sensibilité, scénarios de stress.

\subsection{Page ``Alertes \& Décisions''}

\begin{figure}[H]
  \centering
  \includegraphics[width=0.95\textwidth]{page_decision_alerte.png}
  \caption{Alertes métier et recommandations d'action}
  \label{fig:agriyield-vue-alertes}
\end{figure}

\textbf{Fonction~:} traduire la variation de rendement en alertes opérationnelles (danger/vigilance/opportunité) et recommandations priorisées.

\subsection{Page ``Performances Modèles''}

\begin{figure}[H]
  \centering
  \includegraphics[width=0.95\textwidth]{page_analyse_perfomance.png}
  \caption{Importance des variables pour le modèle actif}
  \label{fig:agriyield-vue-performance}
\end{figure}

\textbf{Fonction~:} rendre le modèle interprétable via la contribution relative des variables.

\subsection{Page ``Gestion Modèles''}

\begin{figure}[H]
  \centering
  \includegraphics[width=0.95\textwidth]{page_gestion_modele.png}
  \caption{Catalogue des versions de modèles et gouvernance}
  \label{fig:agriyield-vue-gouvernance}
\end{figure}

\textbf{Fonction~:} renommer, activer, supprimer une version (hors baseline protégée), et revenir au modèle de base.

%--------------------------------------------------
\section{Pourquoi Random Forest apparaît meilleur que XGBoost selon la baseline}

\subsection{État réel du registre de modèles}

La confusion observée est normale et vient de la coexistence de plusieurs versions \textit{baseline}. Le registre distingue~:
\begin{itemize}[leftmargin=1.5cm]
  \item le \textbf{modèle actif} (utilisé en production UI),
  \item le \textbf{baseline canonique} (pointeur de référence),
  \item d'autres baselines historiques conservées pour audit/rollback.
\end{itemize}

\begin{table}[H]
\centering
\caption{Comparaison des versions baseline présentes dans le registre}
\label{tab:baseline-versions}
\pgfplotstabletypeset[
  col sep=comma,
  columns={VersionShort,BestModel,R2,EvalStrategy,Status},
  columns/VersionShort/.style={string type,column name=Version,column type=l},
  columns/BestModel/.style={string type,column name=Meilleur algorithme,column type=l},
  columns/R2/.style={column name=$R^2$,column type=r,fixed,precision=4},
  columns/EvalStrategy/.style={string type,column name=Stratégie évaluation,column type=l},
  columns/Status/.style={string type,column name=Statut,column type=l},
  every head row/.style={before row={\begin{tabular}{|l|l|r|l|l|}\hline\rowcolor{TableHeader}},after row={\\\hline}},
  every last row/.style={after row={\\\hline\end{tabular}}},
]{data_report/baseline_versions.csv}
\end{table}

\textbf{Interprétation~:}
\begin{itemize}[leftmargin=1.5cm]
  \item la version baseline active dans ton interface est \textbf{Random Forest} avec $R^2\approx0.9577$,
  \item une autre baseline plus récente est \textbf{XGBoost} avec $R^2\approx0.8926$, évaluée en \texttt{temporal\_holdout}.
\end{itemize}

\begin{figure}[H]
\centering
\begin{tikzpicture}
\begin{axis}[
  width=0.85\textwidth,
  height=6.6cm,
  ybar,
  bar width=18pt,
  ymin=0.80,
  ymax=1.00,
  ylabel={$R^2$},
  xlabel={Versions baseline},
  symbolic x coords={model-202604...7ab3f2,model-202604...f31275},
  xtick=data,
  x tick label style={rotate=12,anchor=east},
  nodes near coords,
  nodes near coords align={vertical},
  grid=major,
]
\addplot+[fill=UDSBlue!70] table[col sep=comma, x=VersionShort, y=R2] {data_report/baseline_versions.csv};
\end{axis}
\end{tikzpicture}
\caption{Niveaux de performance ($R^2$) des versions baseline}
\label{fig:baseline-versions-r2}
\end{figure}

%--------------------------------------------------
\section{Intérêt de comparer plusieurs algorithmes}

\subsection{Comparaison intra-baseline (algorithmes candidats)}

\begin{table}[H]
\centering
\caption{Performance des algorithmes sur la baseline active/référence du rapport}
\label{tab:baseline-model-metrics}
\pgfplotstabletypeset[
  col sep=comma,
  columns={Model,MAE,RMSE,R2},
  columns/Model/.style={string type,column name=Modèle,column type=l},
  columns/MAE/.style={column name=MAE (hg/ha),column type=r,fixed,precision=0},
  columns/RMSE/.style={column name=RMSE (hg/ha),column type=r,fixed,precision=0},
  columns/R2/.style={column name=$R^2$,column type=r,fixed,precision=4},
  every head row/.style={before row={\begin{tabular}{|l|r|r|r|}\hline\rowcolor{TableHeader}},after row={\\\hline}},
  every last row/.style={after row={\\\hline\end{tabular}}},
]{data_report/baseline_model_metrics.csv}
\end{table}

\begin{figure}[H]
\centering
\begin{tikzpicture}
\begin{axis}[
  width=0.90\textwidth,
  height=6.8cm,
  ybar,
  bar width=14pt,
  ymin=0.80,
  ymax=1.00,
  ylabel={$R^2$},
  xlabel={Algorithmes candidats},
  symbolic x coords={Random Forest,XGBoost,Gradient Boosting},
  xtick=data,
  x tick label style={rotate=8,anchor=east},
  nodes near coords,
  nodes near coords align={vertical},
  grid=major,
]
\addplot+[fill=UDSGold!85] table[col sep=comma, x=Model, y=R2] {data_report/baseline_model_metrics.csv};
\end{axis}
\end{tikzpicture}
\caption{Comparaison des algorithmes sur la même baseline}
\label{fig:baseline-models-r2}
\end{figure}

\textbf{Pourquoi garder plusieurs modèles est utile~:}
\begin{itemize}[leftmargin=1.5cm]
  \item \textbf{robustesse}~: un algorithme peut dominer sur un contexte et perdre sur un autre,
  \item \textbf{contrôle de dérive}~: comparer les écarts au fil des versions détecte plus vite les régressions,
  \item \textbf{compromis métier}~: certains modèles sont plus précis, d'autres plus stables/interprétables,
  \item \textbf{gouvernance}~: la présence de plusieurs candidats permet des décisions d'activation argumentées.
\end{itemize}

%--------------------------------------------------
\section{Prétraitement et analyses de données (modèle de base)}

\subsection{Couverture du dataset}

Le dataset de base utilisé pour l'entraînement couvre~:
\begin{itemize}[leftmargin=1.5cm]
  \item \textbf{45\,729 observations},
  \item \textbf{226 pays},
  \item \textbf{10 cultures},
  \item période \textbf{1990--2024}.
\end{itemize}

\subsection{Courbe de volume annuel et fenêtre de test temporelle}

\begin{figure}[H]
\centering
\begin{tikzpicture}
\begin{axis}[
  width=0.95\textwidth,
  height=7cm,
  xlabel={Année},
  ylabel={Nombre d'observations},
  ymin=0,
  ymax=1400,
  grid=major,
  legend style={at={(0.02,0.98)},anchor=north west},
]
\addplot[draw=none, fill=UDSGold!18] coordinates {(2018,0) (2024,0) (2024,1400) (2018,1400)} \closedcycle;
\addplot[color=UDSBlue, mark=*, mark size=1.2pt, line width=0.9pt]
  table[col sep=comma, x=Year, y=Observations] {data_report/baseline_yearly_analysis.csv};
\legend{Fenêtre test temporelle,Observations/an}
\end{axis}
\end{tikzpicture}
\caption{Distribution annuelle des observations et zone test 2018--2024}
\label{fig:baseline-observations-year}
\end{figure}

\textbf{Lecture~:} les volumes par année restent relativement stables, ce qui rend la comparaison temporelle plus fiable.

\subsection{Évolution du rendement moyen global}

\begin{figure}[H]
\centering
\begin{tikzpicture}
\begin{axis}[
  width=0.95\textwidth,
  height=7cm,
  xlabel={Année},
  ylabel={Rendement moyen (hg/ha)},
  ymin=55000,
  ymax=90000,
  grid=major,
]
\addplot[color=UDSDark, mark=none, line width=1.1pt]
  table[col sep=comma, x=Year, y=MeanYield] {data_report/baseline_yearly_analysis.csv};
\end{axis}
\end{tikzpicture}
\caption{Tendance annuelle du rendement moyen global}
\label{fig:baseline-yield-year}
\end{figure}

\textbf{Lecture~:} la tendance de fond est croissante, avec des fluctuations annuelles qui justifient l'usage de modèles non linéaires.

\subsection{Courbe de prétraitement~: effet de la standardisation}

\begin{figure}[H]
\centering
\begin{minipage}{0.49\textwidth}
\centering
\begin{tikzpicture}
\begin{axis}[
  width=\textwidth,
  height=6.4cm,
  ybar,
  bar width=10pt,
  ymode=log,
  ymin=1,
  ymax=1000000,
  grid=major,
  xlabel={Variable},
  ylabel={Avant standardisation (log)},
  symbolic x coords={Year,Rain,Pesticides,Temp},
  xtick=data,
  nodes near coords,
  nodes near coords align={vertical},
]
\addplot+[fill=UDSBlue!70] table[col sep=comma, x=Feature, y=RawStd] {data_report/baseline_scaling_plot.csv};
\end{axis}
\end{tikzpicture}
\end{minipage}
\hfill
\begin{minipage}{0.49\textwidth}
\centering
\begin{tikzpicture}
\begin{axis}[
  width=\textwidth,
  height=6.4cm,
  ybar,
  bar width=10pt,
  ymin=0.95,
  ymax=1.05,
  grid=major,
  xlabel={Variable},
  ylabel={Après standardisation (zoom)},
  symbolic x coords={Year,Rain,Pesticides,Temp},
  xtick=data,
  nodes near coords,
  every node near coord/.append style={
    font=\scriptsize,
    /pgf/number format/fixed,
    /pgf/number format/precision=2
  },
]
\addplot+[fill=UDSGold!80] table[col sep=comma, x=Feature, y=ZStd] {data_report/baseline_scaling_plot.csv};
\end{axis}
\end{tikzpicture}
\end{minipage}
\caption{Impact de la standardisation Z-score~: distribution brute (gauche) et valeurs standardisées (droite)}
\label{fig:baseline-zscore-impact}
\end{figure}

\textbf{Lecture~:} à gauche, les écarts-types bruts sont très hétérogènes. À droite, après Z-score, ils convergent autour de 1, ce qui confirme l'effet attendu de la standardisation et améliore la stabilité de l'entraînement.

%--------------------------------------------------
\section{Fine-tuning~: clarification opérationnelle}

\subsection{Mécanisme implémenté}

Le fine-tuning est déclenché via l'upload CSV (endpoint \texttt{/api/ml/train/upload}). Deux modes sont possibles~:
\begin{itemize}[leftmargin=1.5cm]
  \item \texttt{replace\_dataset=true}~: entraînement uniquement sur le dataset chargé.
  \item \texttt{replace\_dataset=false}~: fusion dataset actif + dataset chargé, puis déduplication.
\end{itemize}

\subsection{Règle de couverture temporelle (point important)}

\begin{itemize}[leftmargin=1.5cm]
  \item Si le modèle actif couvre jusqu'à 2024, son horizon est 2024.
  \item Si un fine-tuning est fait avec un dataset allant jusqu'à 2020 et \texttt{replace\_dataset=true}, le nouveau modèle actif est borné à 2020.
  \item Si \texttt{replace\_dataset=false}, la borne temporelle dépend de l'union des deux datasets.
\end{itemize}

C'est cette logique qui explique les écarts de couverture et parfois les écarts d'algorithme gagnant entre versions.

%--------------------------------------------------
\section{Conclusion}

Le chapitre montre une ligne complète et cohérente~: architecture, écrans, gouvernance des versions et preuves quantitatives. La confusion ``Random Forest vs XGBoost'' est levée en distinguant clairement \textbf{version active}, \textbf{baseline canonique} et \textbf{baseline de meilleure performance}. La comparaison multi-modèles devient ainsi un atout méthodologique, et non une ambiguïté.

\newpage
